\begin{textsample}{POS Dim 4 – pi – Score 14.00 – pi/pi\_inf\_e\_ef\_4.txt}  \label{ex:f4_pos_041}
1.
CENÁRIO EDUCACIONAL PIAUIENSE : EDUCAÇÃO, CONTEXTOS E DIÁLOGOS Neste capítulo, evidenciam -se alguns aspectos do \textbf{território} piauiense que são essenciais para a compreensão do cenário em que o presente currículo está inserido. 1.1 O Piauí na História : passado e presente Localizado a noroeste da região \textbf{Nordeste}, o Estado do Piauí dispõe de uma área de 251.611 km², o que representa 2,95 \% do \textbf{território} brasileiro, ocupando a 18ª posição no ranking populacional entre as 27 unidades da federação, onde vivem 3.118.360 piauienses de carne, osso e sonhos.
Embora os primeiros movimentos colonizadores realizados pelos portugueses remontem ao século XVII, quando bandeirantes e entradistas penetraram o \textbf{território} em busca de \textbf{indígenas} para mão de obra ( atividade de preação ), para, logo em seguida, instalar as primeiras fazendas de gado, a relação do homem com nosso \textbf{território} é bem mais remota, recuando há, pelo menos, 50 mil anos, segundo atestam as pesquisas produzidas na serra da Capivara, levadas a efeito pela missão franco-brasileira, liderada pela arqueóloga Niede Guidon, desde 1972.
As escavações produzidas no Sítio Arqueológico Boqueirão da Pedra Furada revolucionaram não somente as pesquisas, mas também aspectos importantes da paisagem social e urbana das comunidades que formam a região.
A economia local se fortaleceu bastante graças ao \textbf{turismo} \textbf{arqueológico} que, além de ser uma ótima fonte de renda, eleva a autoestima de nossa gente, fortalece o \textbf{artesanato} à base de cerâmica e, principalmente, auxilia na divulgação das pesquisas que “ descortinam ” o passado e lançam luzes sobre um futuro de \textbf{esperança}.
Prova desse novo cenário foi o aumento crescente de pesquisadores, visitantes, curiosos e aventureiros que lotam a rede hoteleira da região.
Esses grupos são formados, principalmente, por estudantes da educação básica e também superior que todos os dias visitam o Parque Nacional da Serra da Capivara, onde ficam instalados o Museu do Homem Americano e o da Natureza, em busca de novos saberes sobre nossas origens.
Durante o longo processo de exploração que se materializou na estruturação das primeiras fazendas, povoações, vilas e cidades, um cruento, complexo e nem sempre próspero caminho foi desbravado pelos piauienses.
Nesta jornada, homens e \textbf{mulheres} sempre deram prova de seu valor, de sua vocação para a luta e de sua bravura quando lutavam contra a exploração e desmandos das elites locais.
Mais tarde, ainda no século XIX, em face da necessidade de apoiar o processo de emancipação política da colônia do jugo português, o Piauí, mais uma vez, foi palco de resistência quando protagonizou, na vila de Campo Maior, uma das principais guerras desse longo e nada pacífico processo de ruptura, a Batalha do Jenipapo, às margens do \textbf{riacho} que leva o mesmo nome.
Esse processo de adesão e luta da província do Piauí levou à promoção da Vila de Oeiras ( 1758 ) à condição de primeira \textbf{capital}, situação que duraria até 1852, quando a \textbf{capital} seria transferida para Teresina, na Vila do Poti.
A adesão da província ao processo de Independência prenunciava novos tempos de mudança e desenvolvimento, uma vez que o aparelho burocrático do Estado – principal indutor de tais mudanças –, sempre ausente, agora deveria se fazer mais presente, reduzindo o poder dos proprietários rurais e fortalecendo as relações com o poder central, localizado no Rio de janeiro.
Nesse novo cenário, o poder das famílias apresentava sinais de fragilidade e desgaste e entravam em cena “ novos canais de mobilidade social, pelo menos para os homens ( LEWIN, 1993, p. 173 ) e a expansão das oportunidades, vem, nesse novo momento, por meio da educação, com as faculdades de Direito e de Medicina ofertando “ qualificação formal por vocação e por carreiras profissionais ” ( idem ), cujo papel era “ libertar os filhos dos ricos da terra e afrouxar consideravelmente os laços do controle patriarcal, presentes desde os começos ” ( LEWIN, 1993, p. 173 ).
Nessa perspectiva, a relação entre família e poder, que foi sempre uma variável muito presente no processo de formação social do Piauí, legitimando o prestígio e a ocupação dos mais importantes cargos, postos e funções no aparelho burocrático do Estado, sofre visíveis abalos.
Mudanças significativas nesse cenário foram percebidas principalmente a partir da segunda metade do século XX, quando os primeiros sopros modernizadores começam a ser incorporados ao Estado.
Essa modernização, vista sob a ótica da dimensão política, deveria conferir ao aparelho estatal as condições necessárias para atender às demandas sociais de forma satisfatória e ao mesmo tempo dotar os indivíduos de condições políticas, a fim de que aumentasse de modo qualitativo e plural sua capacidade de participação na vida do Estado.
A década de 70, do século passado, reservou importantes avanços, ainda que de ordem desigual, no sentido de acelerar o desenvolvimento do estado brasileiro e de todas as suas regiões.
Foi nesta década que se iniciou a difusão da ideia de desenvolvimento sustentável, de preservação das espécies e principalmente de que o “ progresso a qualquer custo ” deveria ser substituído por uma solução racional, eficiente e principalmente que não comprometesse as atuais e futuras gerações.
Importante decisão nessa lógica foi a inclusão do planejamento e da modernização econômica, ações que fortaleceriam as demais medidas já em curso.
Assim, o governo estadual implantou, em 2007, um modelo sustentável e eficiente que objetiva dotar de racionalidade e agilidade o Estado e que representou um novo paradigma na governança colaborativa.
Trata -se de um modelo que preconizou uma nova conformação administrativa do espaço geográfico do estado, reorganizando -o em 28 aglomerados, 4 Macrorregiões e 11 \textbf{Territórios} de desenvolvimento.
Esse modelo de governança e gestão foi estruturado considerando as características ambientais, vocações produtivas e dinamismo das regiões, além das relações socioeconômicas e culturais estabelecidas entre as cidades de cada \textbf{território} ; regionalização político-administrativa e melhor malha viária existente.
Tal modelo tinha como propósito reduzir as desigualdades e propiciar a melhoria da qualidade de vida da \textbf{população} piauiense, através da democratização dos programas e ações da regionalização do orçamento.
Em dez anos, segundo dados do IBGE, o Piauí progrediu muito nas diferentes áreas, mas foi na educação que se observou um movimento mais importante, alcançando e superando a meta proposta pelo estado para o IDEB tanto nos Anos Iniciais quanto nos Finais do Ensino Fundamental.
Em relação ao Ensino Médio, o Piauí deixou um desconfortável 24º para o 16º lugar no ranking global do Brasil.
Em 2007, segundo este mesmo levantamento, o Piauí figurava na 7ª posição entre os estados da região \textbf{Nordeste}, saltando em 2017 para a 4ª posição.
Ainda em relação a esse nível de escolaridade, a meta é avançar no número de matrículas, reduzir drasticamente a evasão e aumentar de forma progressiva e sustentável o número de escolas com resultados acima das metas estabelecidas pelo MEC, atingindo no menor tempo possível o topo dessa pirâmide, superando estados como Ceará e Pernambuco, atualmente com os melhores índices.
Nessa perspectiva, os esforços estão orientados no sentido de fazer crescer nosso Índice de Desenvolvimento Humano ( IDH ) e, para tanto, o Piauí está desenvolvendo ações que aumentem o acesso, a permanência e o aprendizado de quem precisa e deve estudar.
Assim, a meta é a construção de um Currículo que respeite a diversidade e assegure o direito às diferenças dos alunos, em suas distintas realidades e em consonância com as orientações da BNCC, homologada em 2017.
Com essa combinação de ações, espera -se alcançar a tão sonhada educação de qualidade, garantindo também uma formação voltada para o mundo do trabalho e para os direitos humanos, aumentando as vagas na Educação de Jovens e Adultos, flexibilizando o currículo, diversificando o ensino, para também atender aos povos \textbf{indígenas}, quilombolas, \textbf{ciganos} e às pessoas com deficiência e transtornos globais do desenvolvimento, a fim de garantir a todos o direito de aprender.

% matched lemmas: arqueológico, artesanato, capital, cigano, esperança, indígena, mulher, nordeste, população, rio, território, turismo
\end{textsample}
